% ECONOMICS_PROBLEM: {"id": "C79-13", "title": "Determinacy of Dadaist Delta games", "jel_code": "C79", "source_id": "OP-0219"}
% !TeX program = pdflatex
\documentclass[11pt,a4paper]{article}
\usepackage[margin=22mm]{geometry}
\usepackage[T1]{fontenc}
\usepackage[utf8]{inputenc}
\usepackage{lmodern,amsmath,amssymb,amsthm,mathtools}
\usepackage{booktabs,longtable,array,enumitem,xcolor,xurl,hyperref,listings}
\hypersetup{colorlinks=true,linkcolor=blue,urlcolor=blue}
\setlength{\parindent}{0pt}
\setlength{\parskip}{5pt}
\setlength{\emergencystretch}{3em}
\newtheorem{theorem}{Theorem}[section]
\newtheorem{lemma}[theorem]{Lemma}
\newtheorem{proposition}[theorem]{Proposition}
\newtheorem{corollary}[theorem]{Corollary}
\theoremstyle{definition}\newtheorem{definition}[theorem]{Definition}
\theoremstyle{remark}\newtheorem{remark}[theorem]{Remark}
\newcommand{\R}{\mathbb R}\newcommand{\Q}{\mathbb Q}
\newcommand{\N}{\mathbb N}\newcommand{\E}{\mathbb E}
\newcommand{\1}{\mathbf 1}
\DeclareMathOperator*{\argmax}{arg\,max}
\DeclareMathOperator*{\argmin}{arg\,min}
\newcommand{\supp}{\operatorname{supp}}
\newcommand{\conv}{\operatorname{conv}}
\newcommand{\rank}{\operatorname{rank}}
\lstset{basicstyle=\ttfamily\scriptsize,breaklines=true,columns=fullflexible,
 keepspaces=true,showstringspaces=false,frame=single,aboveskip=8pt,belowskip=8pt}
\begin{document}
\begin{center}
{\Large\bfseries Determinacy of Dadaist Delta-sums}\par
\medskip
{\large C79-13: research attempt and adversarial verification}\par
9 October 2026
\end{center}
\textbf{Tools.} Web browsing: YES. Exact rational computation: YES.\par
\section{FORMAL RESTATEMENT}
Work in ZFC. Every component $G_i$, indexed by a set $I$, is a fixed set-sized combinatorial game form. Each follower has a set of Left options and a set of Right options. The union of these option relations is well founded, meaning it has no infinite descending sequence; forms, not Conway-equivalence classes, are the data.

A move in $\mathbf G=\sum_{i\in I}^{\Delta}G_i$ changes exactly one coordinate to an option for the moving player. At a limit ordinal, each coordinate is replaced by its eventual value. A fixed player, called First, begins after every limit and at time zero; within a finite block of successive moves the players alternate. A player who is due to move and has no legal move loses. A strategy is a legal move prescription on every nonterminal history where its player is due to move. It is winning when every maximal compatible play is won by that player.

The target is
\[
 \forall I\ \forall(G_i)_{i\in I}\ \forall a\in\{\mathrm L,\mathrm R\}:\quad
 W_{\mathrm L}(\mathbf G,a)\lor W_{\mathrm R}(\mathbf G,a).
\]
The attachment's sentence beginning `Equivalently' asks instead for an undetermined sum. That existential statement is the logical negation of the displayed universal statement, not an equivalent assertion. The minimal correction is to ask for a proof of the universal statement or a counterexample to it.

The countable subproblem restricts $I$ to $\N_0=\{0,1,\ldots\}$. There is no draw convention at $\omega$. Empty index sets and families with only zero components are allowed and are immediate losses for First. Neither an a priori finite play-length bound nor finite branching is assumed.

\textbf{Stress points.} Termination at some ordinal does not by itself justify backward induction; strategies must work against all opponents' plays; a sum with infinitely many coordinates can have an infinite legal run even though no single coordinate can.

\section{QUICK LITERATURE/CONTEXT CHECK}
Paolo Lipparini, \emph{Infinite sums of combinatorial games (Dadaist games)}, arXiv:2505.00614v1, Definitions 3.1 and 3.3 and Problem 3.4, asks exactly the determinacy question. Example 3.6 already gives the infinite impartial case, and Fact 3.5 covers finite sums. \url{https://arxiv.org/html/2505.00614v1}.

The work below supplies full proofs of the foundational bounds and extends the impartial exhaustion argument to the explicitly defined dicot subclass. It does not infer determinacy of unrestricted partizan sums from the source's examples. The inspected version's open-problem designation is context, not a mathematical obstruction to trying further constructions.

\section{ATTACK PLAN}
\textbf{Proof strategies.} First test well-founded recursion on the whole sum; next try an exhaustion strategy using a well-order of active coordinates; finally ask whether an ordinal-valued rank survives limit stages.

\textbf{Construction strategies.} Seek a partizan example in which one player can leave resources inaccessible to the other; test unilateral ends; investigate whether terminal availability after a limit can encode an undetermined infinite-choice payoff. The exhaustion strategy succeeds on a broad subclass, while the general encoding/rank step remains missing.

\textbf{Tools and their roles.} Well-foundedness bounds moves per coordinate; Choice well-orders active coordinates and selects legal options; cardinal arithmetic bounds total play length; limit-stage stabilization validates positions; well-founded recursion proves finite-support determinacy; priority scheduling exhausts dicot resources; impartial finite-game recursion tests boundary cases; unilateral ends diagnose failure of the scheduling argument.

\textbf{Tiny cases.} $0$ is a loss for First; $*=\{0\mid0\}$ is a win for First; two copies of $*$ are a loss for First; three are a win for First. One-sided $1=\{0\mid\}$ is won by Left whichever player starts. Countably infinitely many copies of $*$ are won by Second under the limit rule, not drawn.

\section{WORK}
\begin{lemma}[Finite coordinate activity and a set-sized play bound]\label{lem:bound}
In every legal run, each coordinate changes only finitely many times. Limit positions therefore exist. If $I$ is infinite with cardinality $\kappa$, every run has length less than the successor cardinal $\kappa^+$. For finite $I$, every run has finite length.
\end{lemma}
\begin{proof}
If one coordinate changes infinitely often, select its first $\omega$ changes in their ordinal order. The associated positions form an infinite descending option sequence, contradicting well-foundedness. Each coordinate is consequently constant after its last change before any specified limit, which defines its eventual value.

For infinite $I$, associate each move with the pair consisting of its coordinate and the positive integer counting how many times that coordinate has changed so far. This association is injective. Thus the set of move times injects into $I\times\N$, a set of cardinality $\kappa$ in ZFC. A run with length at least the initial ordinal $\kappa^+$ would give an injection of $\kappa^+$ into a set of cardinality $\kappa$, impossible. For finite $I$, an infinite sequence of moves would change at least one of the finitely many coordinates infinitely often, again a contradiction.
\end{proof}

\begin{remark}
For countably many coordinates the lemma bounds each run by a countable ordinal, not by $\omega$ and not by a uniform finite integer. For example, with countably many copies of $*$ the players may jointly exhaust one infinite subfamily in the first $\omega$ moves while leaving another infinite subfamily untouched, so a legal run can continue beyond $\omega$. The second player's winning strategy below does not require the opponent to cooperate with such a schedule.
\end{remark}

\begin{lemma}[Finite-support determinacy]\label{lem:finite}
If only finitely many components are nonterminal initially, the game is determined for each choice of First.
\end{lemma}
\begin{proof}
Terminal components can never become nonterminal and can be deleted. By Lemma~\ref{lem:bound} the remaining product has no infinite move sequence. The set of positions and the player-to-move indicator therefore has a well-founded successor relation: otherwise dependent choice, available in ZFC, would select an infinite descending sequence.

Define a position to be winning for its player to move if it has a successor losing for the next player; define it to be losing otherwise. This definition is legitimate by well-founded recursion. At a winning position, choose one losing successor. At a losing position, every legal successor is winning for the opponent; a position with no legal successor is losing. These prescriptions prove by well-founded induction that exactly the appropriate player has a winning strategy. Choice selects successors simultaneously when necessary. No limit move is reached in these finite-support plays.
\end{proof}

\begin{definition}[Dicot form]
A form is dicot if every follower, including the initial form, is either terminal for both players or has at least one option for each player. The two option sets need not be equal. Every impartial form is dicot, but the definition also permits partizan option sets.
\end{definition}

\begin{theorem}[Infinite dicot sums are second-player wins]\label{thm:dicot}
Suppose every component is dicot and infinitely many are nonterminal. Let the number of initially nonterminal components have infinite cardinality $\kappa$. For either choice of First, Second has a strategy under which play ends after exactly $\kappa$ moves and First loses at that limit stage. In particular the game has outcome $\mathrm P$.
\end{theorem}
\begin{proof}
Discard initially terminal coordinates and use Choice to index the others by the initial ordinal $\kappa$. Second always chooses the least-index nonterminal coordinate and any legal option there for Second. Dicotness makes that prescription legal whenever any coordinate is nonterminal; fix an option selection on the set of all eligible followers in advance.

Before any stage $\lambda<\kappa$, fewer than $\kappa$ coordinates have been touched, since at most one is touched per move and $\kappa$ is an initial ordinal. An untouched, initially nonterminal coordinate therefore remains. Dicotness implies both players have a legal move somewhere. Thus the run cannot end before stage $\kappa$, regardless of First's strategy.

There are $\kappa$ Second-move times below $\kappa$. One way to verify this cardinal assertion is to write each ordinal as a limit ordinal (or zero) plus a finite integer and inject the finite offsets into the even positive offsets in the same block. Such offsets remain below an infinite initial ordinal. For $\kappa=\omega$, these are simply the even-numbered moves.

Suppose coordinate $\xi<\kappa$ were still nonterminal at stage $\kappa$. Since terminal coordinates cannot change again, it was nonterminal at every earlier stage. Every Second move before $\kappa$ would then have been made in a coordinate of index at most $\xi$, by the least-index rule. Each such coordinate changes only finitely often by Lemma~\ref{lem:bound}. If $\kappa=\omega$, this is a finite set of coordinates with a finite total number of changes. If $\kappa>\omega$, its total number of changes is at most
$\max(|\xi+1|,\aleph_0)<\kappa$. Both conclusions contradict the $\kappa$ Second moves. Therefore every coordinate is terminal at stage $\kappa$.

The ordinal $\kappa$ is a limit, and First is designated to move there by the rule in the problem. There is no legal move. First loses. The strategy and argument did not distinguish Left from Right, proving the claim for both starting choices.
\end{proof}

\begin{corollary}
Every set-indexed dicot $\Delta$-sum is determined: finite nonterminal support is covered by Lemma~\ref{lem:finite}, and infinite nonterminal support by Theorem~\ref{thm:dicot}.
\end{corollary}

\subsection{Why a general proof has not been obtained}
The cardinal bound is not a well-founded rank of the successor relation. For example, the positions obtained by successively removing finitely many stars from an infinite star sum form an infinite chain of legal successors. A rank decreasing strictly at every move would forbid that chain and thus cannot be the naive finite-sum rank.

The exhaustion strategy is not available on general partizan forms. If all active components are $1=\{0\mid\}$ and Right is Second, Right has no legal move on the first response. More generally, First can move a component to a follower in which only First has options. Dicotness was used to exclude precisely this phenomenon. The one-sided example is determined and is not a counterexample to the original conjecture; it only breaks the proposed general scheduling proof.

\section{VERIFICATION}
\textbf{Quantifiers and set theory.} The strategy in Theorem~\ref{thm:dicot} is a single prescription and works against every First strategy. It uses the cardinality of initial nonterminal support, not an arbitrary well-order type whose proper initial segments might already have that cardinality. This distinction matters for singular cardinals. The proof uses only that each proper initial segment of an initial ordinal has smaller cardinality; it does not assume regularity of $\kappa$. Limit stabilization and termination are proved separately. Both players cannot have winning strategies: running the two prescriptions against each other gives one terminating play and one winner.

\textbf{Computational scope.} No finite program tests all transfinite partizan games. The following exact finite check only verifies the stated star boundary cases and the finite recursion convention; it is not used to infer the infinite theorem.
\begin{lstlisting}[language=Python]
from functools import lru_cache
@lru_cache(None)
def win_stars(k):
    # All k moves lead to the same (k-1)-star position.
    if k==0: return False
    return not win_stars(k-1)
assert [win_stars(k) for k in range(7)] == [False,True,False,True,False,True,False]
print('finite star checks: 0 through 6 verified')
# Finite-game generalization:
# Win(position,player) = any(not Win(successor,other_player)
#                            for successor in legal_successors)
# This recursion is used only when the successor relation is well founded.
\end{lstlisting}

\section{FINAL}
\textbf{LABEL: UNRESOLVED for arbitrary partizan $\Delta$-sums.}

\textbf{Strongest fully proved partial result.} All set-indexed dicot sums are determined; if infinitely many components are nonterminal, Second wins by priority exhaustion at their initial cardinal. The impartial specialization is already present in the cited source. The proof also establishes the play-length bound and finite-support determinacy without a finite-branching hypothesis.

\textbf{Exact first gap.} Determinacy has not been proved for a family with infinitely many active coordinates and followers that offer moves to only one player. No reduction of these one-sided residual resources to a determined game, and no undetermined example in this class, has been established here.

\textbf{Three next targets.} First, classify sums of well-founded one-sided ends together with a dicot remainder under the reset-at-limits convention. Second, seek a strategy-invariant comparison of the two players' accessible residual coordinates that survives a limit. Third, test an explicit encoding of an undetermined length-$\omega$ win set using only independent well-founded component forms; the missing requirement is to realize its payoff through legal post-limit move availability.

\textbf{Minimal-counterexample information.} Any counterexample must have infinitely many initially nonterminal coordinates and at least one reachable one-sided follower; it cannot belong to the dicot class. This argument gives no lower bound above countable cardinality and does not assert the existence of a counterexample.

\end{document}
